\documentclass[verbale=interno]{gpt}

\titolo{Struttura del repository della documentazione}
\data{2026-10-20}
\orario{21:00}{22:30}
\luogo{Discord}
\redattore{guglielmo}
\verificatore{corro}
\componenti{biasini,cescobolla,corro,guglielmo,mazzurana,spagnoli}
\assenti{biasini,cescobolla,corro,mazzurana}

\modifica{1.0.0}{2026-10-14}{cescobolla}{corro}{Verifica e approvazione}
\modifica{0.1.0}{2026-10-22}{guglielmo}{--}{Prima stesura}

\begin{document}

\section{Ordine del giorno}
\begin{enumerate}
  \item Presentazione della struttura proposta per il repository della
    documentazione.
  \item Discussione e modifiche alla struttura.
  \item Prossimi passi per il caricamento nell'organizzazione GitHub.
\end{enumerate}

\section{Svolgimento}

La riunione è la call straordinaria prevista nel verbale del 19 ottobre
2026 per avviare l'organizzazione dei repository. Vi hanno partecipato
\nomi{guglielmo,spagnoli}.

\subsection{Presentazione della struttura proposta}
\nome{guglielmo} ha presentato la struttura che aveva preparato per il
repository \texttt{docs}: il template \LaTeX{} comune a verbali e
documenti, la compilazione automatica dei PDF e il sito su cui vengono
pubblicati. Nella proposta i file erano divisi per tipo (verbali,
documenti, lettere), mentre la suddivisione per milestone compariva solo
sul sito.

\subsection{Modifiche alla struttura}
I partecipanti, su suggerimento di \nome{spagnoli}, hanno deciso di cambiare la struttura per renderla più
chiara: la suddivisione per milestone (Candidatura, RTB e PB) deve
esserci anche nel repository, non solo sul sito. Ogni milestone ha una
propria cartella, con all'interno i documenti e i verbali, a loro volta
divisi in interni ed esterni. In questo modo chi cerca un file lo trova
nello stesso posto in cui lo trova sul sito.

\subsection{Prossimi passi}
Prima di caricarla nell'organizzazione, la struttura viene provata in un
repository di prova, così da non introdurre errori nel repository
ufficiale.

\begin{nota}[Aggiornamento successivo alla riunione]
Dopo l'incontro \nome{spagnoli} ha fatto notare a \nome{guglielmo} che
la struttura si poteva migliorare ulteriormente. Si è quindi deciso di
lavorarci ancora prima del caricamento: suddividere meglio le cartelle
e migliorare il sito. Fino ad allora il lavoro è proseguito nel
repository di prova.
\end{nota}

\section{Decisioni}
\begin{decisioni}
  \decisione{Il repository della documentazione è suddiviso per
    milestone: Candidatura, RTB e PB. Ognuna contiene documenti e
    verbali, divisi in interni ed esterni. La stessa suddivisione è
    usata sul sito.}
  \decisione{La struttura viene provata in un repository di prova e
    caricata nell'organizzazione solo quando è completa e funzionante.}
\end{decisioni}

\section{Azioni da intraprendere}
\begin{azioni}
  \azione{Riorganizzare il repository di prova secondo la nuova
    struttura e migliorare il sito.}{guglielmo}{26/10/2026}
  \azione{Caricare la struttura nel repository \texttt{docs}
    dell'organizzazione con una pull request.}{guglielmo}{26/10/2026}
  \azione{Presentare la nuova struttura al resto del
    gruppo.}{guglielmo,spagnoli}{26/10/2026}
\end{azioni}

\section{Prossima riunione}
La prossima riunione ordinaria è lunedì 26 ottobre 2026 alle 21:00 su
Discord.

\end{document}
